\begin{table}[!h]
\centering
\caption{\label{tab:R2_large}Variance Decomposition Across Scenarios Under the Large-Confounding Condition}
\centering
\fontsize{10}{12}\selectfont
\begin{threeparttable}
\begin{tabular}[t]{lcccccc}
\toprule
Scenario & $R_C^2(X)$ & $R_C^2(Y)$ & $R_{\text{incr}}^2(X)$ & $R_{\text{incr}}^2(Y)$ & $R_{\text{total}}^2(X)$ & $R_{\text{total}}^2(Y)$\\
\midrule
Scenario 0 & 0.00 & 0.00 & 0.04 & 0.13 & 0.04 & 0.13\\
Scenario 1 & 0.47 & 0.38 & 0.27 & 0.53 & 0.74 & 0.91\\
Scenario 2 & 0.47--0.56 & 0.24--0.38 & 0.27--0.32 & 0.40--0.53 & 0.74--0.89 & 0.65--0.91\\
Scenario 3A & 0.47--0.56 & 0.14--0.38 & 0.27--0.31 & 0.32--0.53 & 0.74--0.88 & 0.46--0.91\\
Scenario 3B & 0.43--0.50 & 0.35--0.42 & 0.27--0.28 & 0.51--0.54 & 0.70--0.78 & 0.86--0.96\\
\addlinespace
Scenario 4A & 0.47--0.51 & 0.32--0.38 & 0.27--0.29 & 0.47--0.53 & 0.74--0.80 & 0.78--0.91\\
Scenario 4B & 0.47 & 0.38--0.39 & 0.27--0.28 & 0.53 & 0.74--0.75 & 0.91--0.92\\
Scenario 5A & 0.47--0.51 & 0.27--0.38 & 0.27--0.29 & 0.43--0.53 & 0.74--0.80 & 0.69--0.91\\
Scenario 5B & 0.47 & 0.38--0.41 & 0.27--0.28 & 0.53--0.55 & 0.74--0.75 & 0.91--0.96\\
Scenario 5C & 0.45--0.48 & 0.37--0.40 & 0.27 & 0.52--0.53 & 0.72--0.76 & 0.88--0.93\\
\addlinespace
Scenario 5D & 0.46--0.47 & 0.38 & 0.27 & 0.53 & 0.74 & 0.91\\
Scenario 6 & 0.47--0.52 & 0.35--0.43 & 0.27--0.29 & 0.51--0.55 & 0.74--0.81 & 0.87--0.97\\
\bottomrule
\end{tabular}
\begin{tablenotes}
\item \textit{Note: } 
\item $R_C^2$ denotes variance explained by baseline confounders. $R_{\text{incr}}^2$ denotes the incremental variance explained by lagged dynamic processes. $R_{\text{total}}^2$ denotes total variance explained. Values represent minimum--maximum ranges across time points $t = 1, \dots, 5$.
\end{tablenotes}
\end{threeparttable}
\end{table}
